\documentclass[11pt]{article}
\usepackage[T1]{fontenc}
\usepackage[utf8]{inputenc}
\usepackage[a4paper,margin=2.4cm]{geometry}
\usepackage{amsmath,amssymb}
\pagestyle{plain}
\setlength{\parskip}{4pt}
\newcommand{\A}{\mathcal A}
\newcommand{\G}{\mathcal G}

\begin{document}

\begin{center}
{\Large\bfseries Erratum (v3): Theorem 3.1 not established;\\
withdrawal of Corollaries 3.2--3.3}\\[5pt]
{\large Geodesic Convexity and Structural Limits of Curvature Methods\\
for the Yang--Mills Mass Gap on the Lattice}\\[6pt]
Lluis Eriksson \quad\textbullet\quad ai.viXra:2602.0036 \quad\textbullet\quad July 2026
\end{center}

\noindent\rule{\textwidth}{0.8pt}

\vspace{2pt}
\noindent
\textbf{Corollaries 3.2 and 3.3 are withdrawn, and Theorem 3.1 is not established by
its printed proof.} The corollaries pass from a Ricci bound at regular points to a
spectral gap for a global operator, by a one-sentence appeal to codimension; that step
is not established, and the reference given for it does not supply it (\S1). But the
regular-point bound itself does not survive either: the O'Neill trace in equation (19)
omits the mixed horizontal--vertical sectional terms, and the omitted terms have the
unfavourable sign (\S2). A correct global statement exists --- but it concerns a
measure and an operator that \emph{have not been identified with those of the
withdrawn corollaries}, and is therefore added here as a separate result rather than
presented as a repair (\S3).

\emph{An earlier draft of this note stated that ``Theorem 3.1 stands exactly as
printed''. That was wrong, and is withdrawn here. It is recorded rather than quietly
amended: it is the fourth time in this correction campaign that a list of surviving
statements was closed by checking dependencies instead of reading a proof.}

\section{What is withdrawn}

Corollary 3.2 concludes
$\lambda_1\ge N_c/2$ for the Laplacian on the compact quotient $B$, arguing that
``\emph{the singular set (reducible connections) has codimension $\ge2$ and does not
affect the spectrum}'', with a single citation. Corollary 3.3 converts that into a
kinetic gap $E_1^{\rm kin}-E_0^{\rm kin}\ge g^2N_c/4$ uniform in $L$.

Three things are wrong with the step, and each alone suffices.

\paragraph{(a) The cited reference is about a different situation.} The citation is to
Bérard-Bergery--Bourguignon, \emph{Laplacians and Riemannian submersions with totally
geodesic fibres}. That is a theorem about \emph{smooth} submersions whose fibres are
totally geodesic. The gauge orbits are not totally geodesic in general, and the
submersion structure degenerates exactly at the reducible connections --- which is
precisely where the statement is needed. It is not a spectral-removability theorem for
a singular stratum.

\paragraph{(b) Codimension $\ge2$ is not by itself a removability criterion.} Two
distinct levels are involved, and they must not be run together. To identify the
Friedrichs/Dirichlet-form realisation across the singular set one would need an
appropriate removability or capacity theorem for the relevant \emph{form} --- that the
closure of the form on $C_c^\infty(B_{\rm reg})$ coincides with the intended global
form. \emph{Essential self-adjointness} of the minimal differential operator, and hence
uniqueness among \emph{all} self-adjoint extensions, is a separate and strictly stronger
question, and is not established here either. Codimension alone gives neither. Moreover
the stabilisers of the reducible configurations in the singular locus have positive
dimension, so the quotient is a stratified space, not merely a manifold with finitely
many quotient singularities.

\paragraph{(c) The paper specifies and justifies no global self-adjoint realisation of
$\mathcal K_B^{\rm reg}$.} $B_{\rm reg}$ is the regular stratum with the singular locus
removed and is in general \emph{incomplete}. The unweighted Laplace--Beltrami
expression on $C_c^\infty(B_{\rm reg})$ is symmetric and non-negative, but it has not
been shown to be essentially self-adjoint; it may admit several self-adjoint
extensions. Realisations \emph{can} be constructed --- the Friedrichs extension of a
chosen form, for one --- so the objection is not that none exists; it is that the paper
picks none and justifies none, while Corollary 3.2 quantifies the spectrum of one.

\section{Theorem 3.1: the printed O'Neill trace is not the right one}

The corollaries would fall on \S1 alone. But the regular-point bound they globalise
does not survive its own proof either.

Equation (19) reads
\[
\operatorname{Ric}_B(X,X)
=\operatorname{Ric}_{\A}(X^h,X^h)
+\tfrac34\sum_i\bigl\|[X^h,E_i^h]^v\bigr\|^2
\;\ge\;\operatorname{Ric}_{\A}(X^h,X^h).
\]
That is not the trace of O'Neill's horizontal equation. With $\{X,E_2,\dots,E_b\}$ an
orthonormal basis of $T_{[U]}B$ and $\{V_1,\dots,V_r\}$ an orthonormal vertical basis,
O'Neill gives $K_B(X,E_i)=K_{\A}(X^h,E_i^h)+\frac34\|[X^h,E_i^h]^v\|^2$, so tracing
over \emph{horizontal} directions only,
\[
\operatorname{Ric}_B(X,X)=\sum_{i=2}^{b}K_{\A}(X^h,E_i^h)
+\tfrac34\sum_{i=2}^{b}\bigl\|[X^h,E_i^h]^v\bigr\|^2 .
\]
The total-space Ricci tensor, however, traces over horizontal \emph{and} vertical
directions:
\[
\operatorname{Ric}_{\A}(X^h,X^h)
=\sum_{i}K_{\A}(X^h,E_i^h)+\sum_{a}K_{\A}(X^h,V_a).
\]
Hence the correct identity is
\[
\boxed{\begin{aligned}
\operatorname{Ric}_B(X,X)={}&\operatorname{Ric}_{\A}(X^h,X^h)
-\sum_{a=1}^{r}K_{\A}(X^h,V_a)\\
&+\tfrac34\sum_{i=2}^{b}\bigl\|[X^h,E_i^h]^v\bigr\|^2
\end{aligned}}
\]
and (19) has dropped $-\sum_a K_{\A}(X^h,V_a)$.

\textbf{The omitted term has the unfavourable sign.} For a bi-invariant metric on a
compact group all sectional curvatures satisfy $K(X,Y)=\frac14\|[X,Y]\|^2\ge0$, so each
$K_{\A}(X^h,V_a)\ge0$ and the omitted term \emph{subtracts}. Whether the O'Neill term
still dominates it is exactly the computation the proof would have to do, and does not.
It cannot be waved through on general grounds: Riemannian submersions do not in general
preserve positive Ricci curvature.

The omitted mixed sectional-curvature trace must be retained \emph{independently} of
whether the fibres are totally geodesic: it arises simply because
$\operatorname{Ric}_{\A}$ traces over vertical directions as well as horizontal ones,
while $\operatorname{Ric}_B$ traces only over horizontal ones. Total geodesicity
controls the O'Neill $T$-tensor, which enters the vertical and mixed equations, not
this horizontal trace. The non-total-geodesicity of the gauge orbits is therefore a
\emph{separate} mismatch --- with the theorem cited in \S1(a) --- and not the reason
these mixed terms occur.

\medskip
\noindent
\textbf{Status.} The statement
$\operatorname{Ric}_B(X,X)\ge\frac{N_c}{2}|X|^2$ on $B_{\rm reg}$ \textbf{is not
refuted}; what is refuted is the printed derivation. \textbf{Theorem 3.1 is not
established by its proof}, and a new argument specific to the gauge quotient would be
required. Until then, \textbf{no pointwise Ricci bound on $B_{\rm reg}$ may be cited
from this paper as available} --- and it is so cited, in
\texttt{ai.viXra:2602.0033} and \texttt{2602.0035}.

\section{The one sound geometric input, on the correct object}

A global result is available, and it does not go through removability at all. With
$\A=SU(N_c)^E$ carrying the product Haar measure, and $\nu=\pi_\#\mathrm{Vol}_{\A}$
the pushforward to $B$, the metric-measure space $(B,d_B,\nu)$ satisfies
$\mathrm{RCD}^*(N_c/4,\dim\A)$, by stability of the curvature-dimension condition under
quotients by compact groups of measure-preserving isometries. This is
ai.viXra:2602.0046 v3, and it handles the singular stratum automatically --- its own
text records that $B$ is generally not a manifold, and that the codimension analysis
is \emph{not needed} for its argument.

That paper also states that its route \emph{bypasses the O'Neill computation}. In the
light of \S2 that remark is worth more than it looked: the synthetic route is now the
\emph{only} sound geometric input in this chain, because the O'Neill route printed
here does not establish its own conclusion.

That result is \textbf{not} a proof of Corollary 3.2. It is a statement about
$K_\nu$, the self-adjoint operator of the closed Dirichlet form on $L^2(B,\nu)$;
Corollary 3.2 was about a realisation of $\mathcal K_B^{\rm reg}$ on
$L^2(B,d\mathrm{Vol}_B)$. On the regular stratum, writing $d\nu=w\,d\mathrm{Vol}_B$
with $w$ the orbit-volume density, the two differential expressions differ by
$\langle\nabla\log w,\nabla\,\cdot\,\rangle$, and they agree only if $w$ is constant
on each relevant component --- which constant stabiliser \emph{type} does not imply,
since the induced metric on the orbits may vary even where all stabilisers are
conjugate.

What makes the $\nu$-side statement the physically natural one is that it is exactly
the gauge-invariant sector: pullback along $\pi$ is an isometry
\[
L^2(B,\nu)\;\cong\;L^2(\A,\mathrm{Vol}_{\A})^{\G},
\]
because $\int_B|f|^2d\nu=\int_{\A}|f\circ\pi|^2 d\mathrm{Vol}_{\A}$ is the definition
of the pushforward. So the replacement for Corollaries 3.2--3.3 should be read as a
statement about gauge-invariant functions on the smooth compact $\A$, where no
singular geometry is involved at all.

\section{Metric normalisation}

$\operatorname{Ric}\ge N_c/2$ for $\langle X,Y\rangle=-\operatorname{tr}(XY)$ and
$\operatorname{Ric}\ge N_c/4$ for $-2\operatorname{tr}(XY)$ are \emph{the same bound in
two normalisations}, not a sharpening of one by the other. This paper's Theorem 3.1
uses the first; ai.viXra:2602.0046 states its condition in the second. Any downstream
use must name its inner product.

The conversion is correct as a formal statement; what is not correct is to read either
side as available. \textbf{Since Theorem 3.1 is not established (\S2), neither the
$N_c/2$ nor the $N_c/4$ pointwise bound may presently be cited from this paper}, and
the v2 parenthetical ``\emph{v1 stated the bound with $N_c/4$; that is true but not
sharp}'' is \textbf{withdrawn}: it belongs to the same defective derivation.

\section{Three citations to a version that does not exist}

This paper cites \texttt{ai.viXra:2602.0035 v2} in three places --- in Remark 3.4, in
Proposition 4.6, and in the discussion of Theorem 5.2 --- for the corrected
ground-state exponent $f=-\log\psi_0^2$ and related material. \textbf{No v2 of
\texttt{2602.0035} exists in the archive.} The file prepared as that version was
mis-filed onto a different record during the replacement campaign of 2026-07-06; the
record \texttt{2602.0035} still shows only its v1. Readers following those citations
will not find what is cited. The three references become live when that record is
corrected; until then they should be read as pointing to material that is prepared but
unpublished.

\section{Scope of this note}

This note audits Section 3 and the object-provenance failure it caused. It is not a
re-audit of the whole paper, and says so rather than leaving the impression of one.

Read for this note and unaffected by \S1--\S2: \textbf{Corollary 2.3} in its corrected
v2 scope, together with the v2 refutation of the v1 claim it replaces; and
\textbf{Theorem 2.5}, whose v1 convexity route is withdrawn in v2 and whose statement
rests instead on the classical Osterwalder--Seiler strong-coupling cluster expansion
--- an external result, correctly cited and not re-derived here. \textbf{Theorem 2.1}
is not affected by the present findings and was not re-audited.

\textbf{Not audited here, and flagged because the paper's headline rests on them:}
Propositions 4.3--4.6 are labelled ``from v1'', were not revised in the v2 pass, and
are stated in this version \emph{without proof}; Remark 4.7 makes the $O(g^2)$
structural ceiling rest on them together with Theorem 4.2. Proposition 4.6 in
particular invokes $\operatorname{Ric}_B$ on the full quotient and therefore inherits
the object-provenance question of \S1--\S2, although it does so in the favourable
direction, as an obstruction rather than a positive bound. Their status is left open
by this note.

\vspace{4pt}
\noindent\rule{\textwidth}{0.8pt}

\noindent\footnotesize
\textbf{Provenance.} The defect of \S1 was found by tracing which paper supplies the
geometric input consumed by the companions, and checking the cited reference rather
than the citing sentence. The distinction that makes it a defect rather than a
citation slip --- a correct theorem about the wrong object for the consuming
dependency --- is not covered by the usual separation of false statement from invalid
proof, and is recorded as its own category on the page that follows. \textbf{The
defect of \S2 came from an external audit of the first draft of this very note}, which
had certified Theorem 3.1 as standing; it was found by re-deriving the O'Neill trace
rather than by reading the surrounding text, and it is the reason \S2 exists. The
page that follows is reproduced verbatim, by instruction: paraphrase is how the
columns drift.

\end{document}
